\documentclass[output=paper,colorlinks,citecolor=brown]{langscibook}
\ChapterDOI{10.5281/zenodo.20611885}
\author{Ursula Doleschal\orcid{}\affiliation{University of Klagenfurt, Department of Slavic Studies and Writing Center}}

\title{The conditions of uninflectability in nouns in the Slavic languages}

\abstract{In Slavic languages, major word classes like nouns, adjectives, and verbs typically inflect for various grammatical categories. However, some lexemes within these parts of speech are uninflected, meaning they do not display any grammatical markers. I first address how to define uninflectedness -- and uninflectability, which is straightforward for nouns but more complex for adjectives and verbs. Next, I explore which lexeme classes are most commonly uninflected in languages such as Czech, Polish, Russian, Serbian, Slovak, and Slovene, focusing primarily on proper names, loanwords, and abbreviations, which are mostly nouns. I highlight that uninflectedness exists on a continuum, with different degrees of uninflectedness within each class, and this continuum is inversely related to other factors like numeral inflection. I argue that the lack of integration of certain nouns into the inflectional system is linked to specific morphological properties like stem ending, base form, and grammatical gender. These properties vary in restrictiveness across languages, with some (e.g., Russian) showing higher uninflectedness sc. uninflectability in nouns, while others (e.g., Slovene) are more permissive. Additionally, noun classes with feminine gender tend to be more restrictive in terms of stem structure and base form than those with masculine gender.}


\IfFileExists{../localcommands.tex}{
   \addbibresource{../localbibliography.bib}
   \usepackage{tabularx,multicol}
\usepackage[hyphens]{url}
\urlstyle{same}

\usepackage{langsci-optional}
\usepackage{langsci-lgr}
\usepackage{langsci-gb4e}

\usepackage{soul}
% \usepackage{booktabs}
% \usepackage{geometry}
\usepackage{multirow}
% \usepackage{makecell} %safely breaks a line inside a table cell
% \usepackage{tablefootnote} %footnotes in table captions
\usepackage{enumerate}
\usepackage{tikz}
\usepackage{array}
\usepackage[shortlabels]{enumitem}

%for having several figures on one line, turning words in 90°angles
% \usepackage{graphicx}
\usepackage{subcaption}

% \usepackage{caption} %for making the caption line longer in concrete cases

% \usepackage{rotating} %rotate a whole page; for large tables

\usepackage{fontspec}
\newfontfamily\charis{Charis SIL} %for being able to use the symbol ‿ and a better command of some of the boldfaced diacritics in 01
\newfontfamily\japanesefont{Noto Serif CJK JP} % for Japanese script
% \usepackage{tipa}

% \usepackage{needspace} %to keep exs on one page

\usepackage{xcolor}

\usepackage{pgfplots}
\pgfplotsset{compat=1.18}

%%%%%   AVMs for 02	%%%%%
\usepackage{langsci-avm}
\avmsetup{switch = ?}


\usepackage{langsci-branding}

   %\newcommand*{\orcid}{}


\makeatletter
\let\thetitle\@title
\let\theauthor\@author
\makeatother

\newcommand{\togglepaper}[1][0]{
   \bibliography{../localbibliography}
   \papernote{\scriptsize\normalfont
     \theauthor.
     \titleTemp.
     To appear in:
     E. Di Tor \& Herr Rausgeberin (ed.).
     Booktitle in localcommands.tex.
     Berlin: Language Science Press. [preliminary page numbering]
   }
   \pagenumbering{roman}
   \setcounter{chapter}{#1}
   \addtocounter{chapter}{-1}
}

\newbool{bookcompile}
\booltrue{bookcompile}
\newcommand{\bookorchapter}[2]{\ifbool{bookcompile}{#1}{#2}}



\renewcommand{\thempfootnote}{\fnsymbol{mpfootnote}} % use * for footnotes in float env

% \newcolumntype{L}[1]{>{\raggedright\arraybackslash}p{#1}} %left-aligned (non-justified) text in tables
% \newcolumntype{R}[1]{>{\raggedleft\arraybackslash}p{#1}} % right-aligned (non-justified) text in tables

\definecolor{customgreen}{RGB}{27,158,119}
\definecolor{custompink}{RGB}{231,42,138}
\definecolor{customblue}{RGB}{117,112,179}
\definecolor{customlightblue}{RGB}{143,170,220}
\definecolor{customdarkblue}{RGB}{68,114,196}
\definecolor{customyellow}{RGB}{225,192,0}
\definecolor{customorange}{RGB}{233,137,21}
\definecolor{customgray}{RGB}{229,229,229}


% Left-aligned indexes
% % \makeatletter
% % \patchcmd{\theindex}{\twocolumn}{\raggedright\twocolumn}{}{}
% % \makeatother

%to make Japanese scripts in the bibliography smaller
\newcommand{\japbib}[1]{{\japanesefont\small #1}}

%%%%%%%%%%%%%%%%%%%	MACROS for 02	%%%%%%%%%%%%%%%%

%%% Formatting	%%%%%

\newcommand{\textex}[1]{\textit{#1}} 		 %for formatting linguistic examples in-text%
%command-shift X

\newcommand{\lxm}[1]{\textsc{#1}}  		%for formatting names of lexemes
%command-option-shift L

%%%%%%%%% Abbreviations	%%%%%%

\newcommand{\featname}[1]{\textsc{#1}}
\newcommand{\featspec}[2]{\featname{#1}:{#2}}
\newcommand{\feat}[2]{\textsc{#1}:{#2}}  				% \feat{FEATURE}{value}  sc's
\newcommand{\featnameonly}[1]{\textsc{#1}} 			% \featnameonly{featurename}  sc's

\newcommand{\pr}{$^{\prime}$}

\newcommand{\Corr}{\textit{Corr}}
\newcommand{\propm}{\textit{pm}}

\newcommand{\PC}{Π$_{C}$}

\newcommand{\PF}{Π$_{F}$}
\newcommand{\PR}{Π$_{R}$}

\newcommand{\lex}{$\mathcal{L}$}

   %% hyphenation points for line breaks
%% Normally, automatic hyphenation in LaTeX is very good
%% If a word is mis-hyphenated, add it to this file
%%
%% add information to TeX file before \begin{document} with:
%% %% hyphenation points for line breaks
%% Normally, automatic hyphenation in LaTeX is very good
%% If a word is mis-hyphenated, add it to this file
%%
%% add information to TeX file before \begin{document} with:
%% %% hyphenation points for line breaks
%% Normally, automatic hyphenation in LaTeX is very good
%% If a word is mis-hyphenated, add it to this file
%%
%% add information to TeX file before \begin{document} with:
%% \include{localhyphenation}
\hyphenation{
    Al-ex-an-der
    Ber-ber
    chan-ges
    cir-cum-stan-ces
    coin-age
    Fraj-zyngier
    ita-liano
    Ja-pa-nese
    Kö-nig
    Krzy-żanowski
    li-te-ra-tur-no-go
    mi-zen-kei
    par-a-digm
    phe-nom-e-non
    ren-tai-shi
    Slo-vene
    Spen-cer
    Stoj-kov
    Ver-bi-ca-re-se
    wide-spread
}

\hyphenation{
    Al-ex-an-der
    Ber-ber
    chan-ges
    cir-cum-stan-ces
    coin-age
    Fraj-zyngier
    ita-liano
    Ja-pa-nese
    Kö-nig
    Krzy-żanowski
    li-te-ra-tur-no-go
    mi-zen-kei
    par-a-digm
    phe-nom-e-non
    ren-tai-shi
    Slo-vene
    Spen-cer
    Stoj-kov
    Ver-bi-ca-re-se
    wide-spread
}

\hyphenation{
    Al-ex-an-der
    Ber-ber
    chan-ges
    cir-cum-stan-ces
    coin-age
    Fraj-zyngier
    ita-liano
    Ja-pa-nese
    Kö-nig
    Krzy-żanowski
    li-te-ra-tur-no-go
    mi-zen-kei
    par-a-digm
    phe-nom-e-non
    ren-tai-shi
    Slo-vene
    Spen-cer
    Stoj-kov
    Ver-bi-ca-re-se
    wide-spread
}

   \boolfalse{bookcompile}
   \togglepaper[23]%%chapternumber
}{}

\begin{document}
\maketitle

\il{Slavic|(}
\section{Introduction} 
In this article I will present some findings of my habilitation thesis \citep{doleschal2000} where I investigated the phenomenon of \isi{uninflectedness} and \isi{uninflectability} in six Slavic languages -- \ili{Russian}, \ili{Polish}, \ili{Slovak}, \ili{Czech}, \ili{Slovene} and \ili{Serbian}. The aim of this investigation was to find out how and why in highly inflecting languages, both single words such as the \ili{Russian} noun \emph{metro} as in \emph{stancija} \emph{metro} `metro station', the \ili{Czech} adjective \emph{blond} in \emph{blond vlasy} `blond hair'or the numeral \emph{dve} in \ili{Serbian} as in \emph{s dve devojke} `with two girls' and entire classes of words\is{inflectional class} as \ili{Russian} and \ili{Polish} nouns ending in the vowel /u/, e.g., \emph{s kenguru} `with (a) kangaroo' remain partly or even fully \isi{uninflected}, and if eventually, interjections as \emph{pryg} `hop' in \ili{Russian} might qualify for \isi{uninflectable} verbs: \emph{my pryg} `we hopped'.

In the course of this research, other phenomena also related to \isi{uninflectedness} emerged, most notably constructions with company names in compound-like structures, such as \emph{Duracell} in \emph{Duracell nagradna igra} (\ili{Slovene}) `Duracell prize draw', \emph{Milka} in \emph{50 Milka plišanih kravica} (\ili{Serbian}) `50 Milka fluffy toy cows' or a historical growth of inflectedness as wittnessed by the (relatively) new \isi{inflectional class} \emph{kuli} in \ili{Czech}, \ili{Slovak} and \ili{Polish}. These issues have been addressed in \citet{doleschal1999, doleschal2008, doleschal2015, doleschal2017}.

The structure of the present article is the following: Firstly, the Slavic languages will be characterised as inflecting languages (Section \ref{03sec2}). Secondly, the different forms of the phenomenon which have been identified in the material will be described (Section \ref{03sec3}). Finally, the ways in which \isi{uninflectability} manifests itself in the six languages under investigation, and how this \isi{uninflectability} is related to the inflectional classes\is{inflectional class} in Slavic, will be examined, as will the research question about the conditions of \isi{uninflectedness} (Section \ref{03sec4}).

\section{Slavic languages as inflecting languages} \label{03sec2}

The Slavic languages are traditionally classified as highly inflecting and belonging to the fusional type (cf. e.g., \citealt[6]{comriecorbett1993-intro}). This means that they signal grammatical meanings by endings which are attached to lexical morphemes, usually stems\is{stem}. Endings are characterised by cumulative exponence, i.e., one ending expresses more than one grammatical meaning (feature value). Inflection is different and characteristic for the different parts of speech.

The inflected parts of speech are: noun, pronoun, adjective, numeral, determiner, verb (including participle; \citealt[222]{sussexcubberley2006}), and -- depending on one's theoretical premises -- adverb (cf. \citealt[743]{nagorko2009}). Within the inflected parts of speech, word forms are organised in paradigms\is{paradigm} forming inflectional classes\is{inflectional class} so that any lexeme of these parts of speech necessarily belongs to (usually) one such class\is{inflectional class} (cf. Table \ref{03tab1}). The grammatical (morphosyntactic and morphosemantic) categories involved in inflection are the following:\footnote{Cf. \citet{spencerwiemer2020} for a discussion of inflectional categories in Slavic languages.} case, number, \isi{gender}, \isi{animacy}, (gradation) in nouns and adjectives; person, number, \isi{gender}, tense, mood, voice, (aspect), gerund, participle in verbs, determining the content \isi{paradigm} of a lexeme that belongs to a certain part of speech. For example, the \ili{Polish} nouns \emph{pan} `gentleman, mister' and \emph{żona} `woman' take different sets of inflectional endings for the same values of the categories case and number (Table~\ref{03tab1}). Furthermore, these two nouns are classified differently with regard to the category of \isi{gender} -- while \emph{pan} has the feature masculine, \emph{żona} displays feminine \isi{gender}. Whether the \isi{gender} feature is also expressed by the sets of endings along with the case and number features is controversial (cf. \citealt[148]{doleschal2009}).

\begin{table}[ht]
	\centering
	\begin{tabularx}{0.8\textwidth}{XXXXX}
		\lsptoprule
			& \textsc{sg} & \textsc{pl} & \textsc{sg} & \textsc{pl} \\
			\midrule
			\textsc{n(om)} & pan-Ø & pan-owie & żon-a & żon-y \\
			\textsc{g(en)} & pan-a & pan-ów & żon-y & żon \\
			\textsc{d(at)} & pan-u & pan-om & żon-ie & żon-om \\
			\textsc{a(cc)} & pan-a & pan-ów & żon-ę & żon-y \\
			\textsc{i(nstr)} & pan-em & pan-ami & żon-ą & żon-ami \\
			\textsc{l(oc)} & pan-u & pan-ach & żon-ie & żon-ach \\
			\textsc{v(oc)} & pan-ie & pan-owie & żon-o & żon-y  \\
		\lspbottomrule
	\end{tabularx}
	\caption{Inflectional classes\is{inflectional class} (\ili{Polish}): \emph{pan} (\textsc{m}) `gentleman, mister', \mbox{\emph{żona} (\textsc{f})} `woman'}
	\label{03tab1}
\end{table}


\section{Types and manifestations of \isi{uninflectedness}} \label{03sec3}

\subsection{General}

In the general area of \isi{uninflectedness} in inflected parts of speech, different manifestations of the phenomenon can be distinguished:

\begin{enumerate}
	\item classes of words\is{inflectional class} that never take any endings (Section \ref{03sec3.2})
	
	\item single words that never take any endings (Section \ref{03sec3.2})
	
	\item words that may or may not take endings and vary by syntactic context, register, sociolect, speaker (or freely) (Section \ref{03sec3.3}).
\end{enumerate}

\subsection{Uninflectability\is{uninflectability} (paradigmatic uninflectedness\is{uninflectedness!paradigmatic})} \label{03sec3.2}

In the first place, there are words that never take any inflectional endings and are therefore paradigmatically uninflected\is{uninflectedness!paradigmatic}, as the borrowings \emph{miss} \textsc{f.} `miss', e.g., in \emph{miss Polonia} or \emph{gnu} \textsc{n.}, \textsc{f.}, \textsc{m.} (in \ili{Russian}, \ili{Polish}, \ili{Slovak}) `gnu' for a type of antelope. Paradigmatic uninflectedness\is{uninflectedness!paradigmatic} pertains to the lexical level. It is an inherent property of the lexeme. This manifestation of \isi{uninflectedness} will be called ``\isi{uninflectability}''.

Uninflectable\is{uninflectable} lexemes can be subdivided into two groups:

\begin{enumerate}
	\item In the first group \isi{uninflectability} is regular, i.e., it can be described by rules based on narrowly linguistic (e.g., phonological or lexical) properties: in \ili{Russian}, e.g., nouns ending in the vowel /u/ cannot be inflected: \emph{tabu} `taboo', \emph{ragu} `ragout', \emph{kenguru} `kangaroo', \emph{Lulu} `Loulou', \emph{Subaru} `Subaru'. This case shall be called \textit{systematic \isi{uninflectability}}.
	
	\item The second group contains words whose \isi{uninflectability} cannot be explained by properties of the language system, as in the case of some toponyms in \ili{Slovak} ending in the vowel /e/ as \emph{Skopje} or in a so-called soft consonant\is{consonant!soft}\footnote{In the Slavic languages, consonants underwent several palatalisations, leading to correlations between soft (palatal or palatalised)\is{consonant!soft} and hard (non-palatal) consonants\is{consonant!hard}. This difference is widely morphologised, so that ``hard'' stems\is{stem} and ``soft'' stems\is{stem} select different sets of endings which are not conditioned by phonology. Moreover, there are ``neutral\is{consonant!neutral}'' sc. ``unpaired'' consonants\is{consonant!unpaired} which do not take part in the correlation (cf. \citealt[6, 9--10]{comriecorbett1993-intro}).\label{03ftn2}} as \emph{Providence} (/s/) while \emph{Opole}, \emph{Bystrice} (both ending in /e/), \emph{Provence} (ending phonologically in /s/) with the same type of phonological ending are inflectable. This case shall be called \textit{unsystematic \isi{uninflectability}}.
\end{enumerate}

On this basis we can now define what exactly will be understood by \textit{\isi{uninflectability}}:

A word is called \isi{uninflectable},\footnote{Cf. also \citet[24--35]{krzyzanowski2013} for a thorough terminological discussion.} if

\begin{enumerate}
	\item it belongs to an inflectable part of speech, and
	
	\item does not take any endings normally found with words of that kind, but
	
	\item may nevertheless appear in any syntactic context typical for the part of speech in question.
\end{enumerate}

Condition 1 is necessary to distinguish the words in question from \isi{uninflected} parts of speech, such as conjunctions and prepositions, in the Slavic languages. Condition 2 distinguishes \isi{uninflectable} words from semi-inflectable ones as, e.g., \ili{Polish} \emph{muzeum} `museum' (Table \ref{03tab2}), which is \isi{uninflected} in the singular but regularly inflected in the plural, and from contextual uninflectedness\is{uninflectedness!contextual} (see below Section \ref{03sec3.3}).

\begin{table}[ht]
	\centering
	\begin{tabularx}{0.4\textwidth}{p{0.5cm}XX}
		\lsptoprule
		& \textsc{sg} & \textsc{pl} \\
		\midrule
		\textsc{n} & muzeum & muze-a \\
		\textsc{g} & muzeum & muze-ów \\
		\textsc{d} & muzeum & muze-om \\
		\textsc{a} & muzeum & muze-a \\
		\textsc{i} & muzeum & muze-ami \\
		\textsc{l} & muzeum & muze-ach \\
		\textsc{v} & muzeum & muze-a \\
		\lspbottomrule
	\end{tabularx}
	\caption{Semi-inflectable nouns in \ili{Polish}}
	\label{03tab2}
\end{table}

Condition 3 is necessary to distinguish between \isi{uninflectability} proper and defectivity (cf. \citealt{doleschal2000, doleschal2002-mici, spencer2020}): for example, the \ili{Russian} indefinite numeral \emph{malo} `a little, few' has only one form and appears solely in Nominative and Accusative Singular expressions (\ref{03ex1}):

\begin{exe}
	\ex \label{03ex1}
	\begin{xlist}
		\ex \gll Tam malo detej. \\
		there few.\textsc{nom} children.\textsc{gen} \\
		\trans ‘There are few children there.’
		
		\ex \gll Ja znaju malo detej. \\
		I know few.\textsc{acc} children.\textsc{gen} \\
		\trans ‘I know few children.’

		\ex \gll *Ona interesuetsja malo detej. \\
		\phantom{*}she is-interested-in few.\textsc{instr} children.\textsc{gen} \\
		\trans ‘She is interested in few children.’
	\end{xlist}
\end{exe}

It is thus also \isi{uninflectable}, but not in the same way as are nouns as \emph{kenguru} `kangaroo' which as such can express all syntactic cases and both singular and plural, e.g. (\ref{03ex2}):

\begin{exe}
	\ex \label{03ex2}
	\begin{xlist}
		\ex \gll Tam molod-oj kenguru. \\
		there young-\textsc{m.nom.sg} kangaroo.\textsc{nom.sg} \\
		\trans ‘There is a young kangaroo there.’
		
		\ex \gll Ja tam vižu kenguru. \\
		I there see kangaroo.\textsc{acc.sg/pl} \\
		\trans ‘I see a kangaroo there.’ / ‘I see kangaroos there.’
		
		\ex \gll Ona interesuetsja kenguru. \\
		she is-interested-in kangaroo.\textsc{instr.sg/pl} \\
		\trans ‘She is interested in a kangaroo / kangaroos.’
	\end{xlist}
\end{exe}

This definition can be exemplified with two female names taken from \ili{Czech} -- \emph{Jana} and \emph{Mimi} (cf. Table \ref{03tab3}). As proper names, both are clearly nouns and thus fulfil condition 1. Secondly, it can be seen from Table \ref{03tab3} that \emph{Mimi} does not take any case endings but occurs with the same case meanings as the inflected lexeme \emph{Jana}, i.e., it has the same content \isi{paradigm} as \emph{Jana}, and thus fulfills condition 2. It can therefore appear in all necessary syntactic contexts as is illustrated by the use of prepositions in the second column and so also fulfils condition 3.

\begin{table}[ht]
	\centering
	\begin{tabularx}{0.45\textwidth}{p{0.5cm}XXX}
		\lsptoprule
		\textsc{n} & & Jan-a & Mimi \\
		\textsc{g} & (bez) & Jan-y & Mimi \\
		\textsc{d} & (k) & Jan-ě & Mimi \\
		\textsc{a} & (pro) & Jan-u & Mimi \\
		\textsc{l} & (při) & Jan-ě & Mimi \\
		\textsc{i} & (s) & Jan-ou & Mimi \\
		\textsc{v} & & Jan-o & Mimi \\
		\lspbottomrule
	\end{tabularx}
	\caption{Inflectable and \isi{uninflectable} name in \ili{Czech}}
	\label{03tab3}
\end{table}

Condition 3 also implies that an \isi{uninflectable} word possesses and represents all morphological categories and their grammatical meanings (feature values) of the part of speech in question. But how large can such a \isi{paradigm} be? Would it also have to cover otherwise periphrastic morphology? Does it have to cover word-class changing morphology?

The \ili{Russian} verb \isi{paradigm} has up to 124 forms (e.g., the verb \emph{ljubit'} `to love'), including the four participles with their paradigms\is{paradigm} and the two gerunds (regarding aspect as a lexical feature, cf. \citealt{genis2021}). Does an \isi{uninflectable} item have to substitute for all of these forms in order to be called \isi{uninflectable} and not defective? For my material, I have come to the following conclusion both on an empirical and a theoretical basis: an \isi{uninflectable} lexeme must be able to independently represent all synthetic forms of the \isi{paradigm} of a part of speech in question. For the nominal parts of speech this means all the case-number forms in the noun and case-number-\isi{gender} forms and adverb in the adjective, and case forms in the numeral as well as case-\isi{gender} forms for some numerals, but not the periphrastic forms, such as periphrastic gradation.

This is mainly due to the fact that words that are usually perceived as \isi{uninflectable} can combine with periphrastic markers such as \ili{Czech} \emph{více} in \emph{více blond} `more blond' for the comparative or as \ili{Russian} \emph{pryg by} `would hop' for the conditional. From the theoretical point of view, it has been argued (cf. e.g., \citealt{haspelmath2000}) that obligatory preposition-noun combinations, such as in the locative in the Slavic languages, are usually seen as phrases although they could just as well be analysed as periphrastic case forms. However, we do not expect an \isi{uninflectable} noun in the (syntactic) locative to co-signal a preposition. Another viewpoint could be that periphrasis is on the borderline between syntax and inflection \citep{brownchumakinacorbettpopovaspencer2012}, hence not a prototypical inflectional marker and could therefore be neglected. As to the word-class changing forms (adverb, participle), \citet{spencer2020} has argued for excluding them from the \isi{paradigm}, and I will follow him here. Under these premises, in \ili{Slovene} an \isi{uninflectable} adjective (as \emph{super}) would substitute for a \isi{paradigm} as presented in Table \ref{03tab4}. 

\begin{table}[ht]
	\centering
	\begin{tabularx}{0.9\textwidth}{p{0.5cm}p{2.6cm}p{1.5cm}p{1.5cm}XXX}
		\lsptoprule
        & \multicolumn{3}{c}{`new'} & \multicolumn{3}{c}{`super'} \\
        \midrule
		& \textsc{m} & \textsc{f} & \textsc{n} & \textsc{m} & \textsc{f} & \textsc{n} \\
		\midrule
		\multicolumn{7}{c}{\textsc{sg}} \\
		\midrule
		\textsc{n} & nov-Ø/nov-i\footnote{Indefinite/definite form (cf. \citealt[36]{greenberg2008}).} & nov-a & nov-o & super & super & super \\
		\textsc{g} & nov-ega & nov-e & nov-ega & super & super & super \\
		\textsc{d} & nov-emu & nov-i & nov-emu & super & super & super \\
		\textsc{a} & nov-Ø/nov-ega\footnote{Inanimate/animate\is{animacy} \isi{gender}.} & nov-o & nov-o & super & super & super \\
		\textsc{l} & nov-em & nov-i & nov-em & super & super & super \\
		\textsc{i} & nov-im & nov-o & nov-im & super & super & super \\
		\midrule
		\multicolumn{7}{c}{\textsc{du}} \\
		\midrule
		\textsc{n} & nov-a & nov-i & nov-i & super & super & super \\
		\textsc{g} & nov-ih & nov-ih & nov-ih & super & super & super \\
		\textsc{d} & nov-im & nov-im & nov-im & super & super & super \\
		\textsc{a} & nov-a & nov-i & nov-i & super & super & super \\
		\textsc{l} & nov-ih & nov-ih & nov-ih & super & super & super \\
		\textsc{i} & nov-imi & nov-imi & nov-imi & super & super & super \\
		\midrule
		\multicolumn{7}{c}{\textsc{pl}} \\
		\midrule
		\textsc{n} & nov-i & nov-e & nov-a & super & super & super \\
		\textsc{g} & nov-ih & nov-ih & nov-ih & super & super & super \\
		\textsc{d} & nov-im & nov-im & nov-im & super & super & super \\
		\textsc{a} & nov-e & nov-e & nov-a & super & super & super \\
		\textsc{l} & nov-ih & nov-ih & nov-ih & super & super & super \\
		\textsc{i} & nov-imi & nov-imi & nov-imi & super & super & super \\
		\lspbottomrule
	\end{tabularx}
	\caption{Inflectable and \isi{uninflectable} adjective in \ili{Slovene}}
	\label{03tab4}
\end{table}

Since gradation can be periphrastic in all Slavic languages, this category need not be expressed by an uninflectable adjective by itself. In \ili{Slovene}, adjectives can form synthetic forms of comparative and superlative as \emph{nov-ejš-i} `new-\textsc{gradation-m.sg}' `newer' and \emph{naj-nov-ejš-i} `\textsc{sup}-new-\textsc{gradation-m.sg}' `newest' along with a periphrastic comparative and superlative, which is then also applied to \isi{uninflectable} adjectives (Table \ref{03tab5}).

\begin{table}[ht]
	\centering
	\begin{tabularx}{\textwidth}{lXXX}
		\lsptoprule
			& Positive & Comparative & Superlative \\
			\midrule
			synthetic & nov-Ø, -a, -o & nov-ejš-i, -a, -e & naj-nov-ejš-i, -a, -e \\
			periphrastic & nov-Ø, -a, -o & bolj nov-i, -a, -o & najbolj nov-i, -a, -o \\
			& super & bolj super & najbolj super \\
		\lspbottomrule
	\end{tabularx}
	\caption{Gradation of adjectives in \ili{Slovene}}
	\label{03tab5}
\end{table}

\subsection{Contextual (syntagmatic) uninflectedness\is{uninflectedness!contextual}} \label{03sec3.3}

Having defined \isi{uninflectability}, let us now turn to a closely related phenomenon which is often treated together with \isi{uninflectability} under the same heading (``nesklonjaemost'', ``neohebnost'') but is distinct from the first: contextual uninflectedness\is{uninflectedness!contextual}.\footnote{In \citet{doleschal2000} “syntagmatic” was used throughout. However, as one of the reviewers pointed out, this term is too narrow. Therefore, the phenomenon is called \textit{contextual uninflectedness}\is{uninflectedness!contextual} here.} Contextual uninflectedness\is{uninflectedness!contextual} means that a word, which is fully inflectable on the lexical level, remains uninflected for contextual\is{uninflectedness!contextual} reasons, such as text type, style, sociolect and last but not least -- syntactic context, e.g., apposition or use of a classifier as in \emph{v naselj-u Njivice} (in village-\textsc{loc.sg} Njivice `in the village Njivice') in \ili{Slovene} or \emph{w województw-ie Białystok} (in province-\textsc{loc.sg} Białystok `in the province Białystok'). However, such \isi{uninflectedness} is usually not obligatory, e.g., in \ili{Polish}, masculine surnames ending in /e/ or /i/ are inflectable, but especially in collocation with a first name or a title often remain \isi{uninflected}, e.g., \emph{Maksimilian-a Kolbe} Maksimilian-\textsc{gen.sg} Kolbe [\textsc{gen.sg}] instead of \emph{Maksimilian-a Kolb-ego} Maksimilian-\textsc{gen.sg} Kolbe-\textsc{gen.sg} (cf., e.g., \citealt[39--40]{krzyzanowski2013}). In \ili{Slovene} this happens with feminine surnames in /a/, e.g., \emph{z Ad-o Muha} `with Ada-\textsc{instr.sg} Muha [\textsc{instr.sg}]' `with Ada Muha' instead of \emph{z Ad-o \mbox{Muh-o}} `with Ada-\textsc{instr.sg} Muha-\textsc{instr.sg}'. In \ili{Serbian}, male surnames can remain \isi{uninflected} in collocation with first names, if they precede the first name: \emph{Petrović Ivan-a} Petrović Ivan-\textsc{gen.sg} `Petrović Ivan's'.

On the other hand, in \ili{Czech} abbreviations as \emph{NATO} or \emph{FIFA} remain \isi{uninflected} in official language, but are inflected in casual speech (\emph{Internetová jazyková příruč\-ka} \citeyear{ujc2024-zkratkovaslova}). Another case in point are \ili{Russian} place names formed with the suffixes \emph{-ovo} and \emph{-ino} such as \emph{Perovo}, \emph{Puškino} which are in principle inflectable and should be inflected according to the norms of Contemporary Standard \ili{Russian} (cf. \citealt{paxomov-nodate, redakcijaportalagramotaru}; \citealt[739]{zaliznjak2008}; \citealt[35]{ageenko2010}). Nevertheless, they often remain \isi{uninflected} both in spoken and written language and seem to show a tendency towards \isi{uninflectability} (\citealt[156--159]{graudina1980}; \citealt[56--59]{panov1968}; \citealt[739]{zaliznjak2008}, but cf. \citealt{tejkin2021} for an opposite view).

One of the interesting issues is if and how contextual\is{uninflectedness!contextual} and paradigmatic uninflectedness\is{uninflectedness!paradigmatic} relate to each other, i.e., is contextual uninflectedness\is{uninflectedness!contextual} a preliminary stage of \isi{uninflectability}? This may be the case as with the aforementioned \ili{Russian} place names, but it is certainly not the case, when an \isi{uninflected} noun occurs together with an appositive classifier in \ili{Slovene} or in \ili{Serbian}. This problem needs an investigation of the relationship between the effects of register and style and language change caused by both sociolinguistic and inherently morphologocial factors, such as productivity of inflectional classes\is{inflectional class}, and has to be left for further exploration.

\subsection{Geographical distribution of \isi{uninflectability} in six Slavic languages}

As a result of the study by \citet{doleschal2000} the manifestation of \isi{uninflectability} in the six Slavic languages under investigation can be characterised as follows: Uninflectability\is{uninflectability} occurs regularly, but in different ways, with the parts of speech \emph{noun}, \emph{adjective}, \emph{(cardinal) numeral}. The degree of \isi{uninflectability} manifests itself as a kind of continuum with two trajectories for different parts of speech: one from west to east and the other from south to north (cf. Table \ref{03tab6}). This means that for nouns the highest degree of \isi{uninflectability} is found in \ili{Russian} and the lowest in \ili{Slovene}, whereas for both adjectives and numerals the highest degree is found in \ili{Serbian} and the lowest (if any) in \ili{Russian}.

\begin{table}[ht]
	\centering
	\begin{tabularx}{0.75\textwidth}{lXXXXX}
		\lsptoprule
			& \textbf{Noun} & & & & \\
			N & Polish\il{Polish} & \cellcolor[RGB]{127,127,127} & \cellcolor[RGB]{51,51,51} & Russian\il{Russian} & N \\
			↑ & Czech\il{Czech} & \cellcolor[RGB]{204,204,204} & \cellcolor[RGB]{178,178,178} & Slovak\il{Slovak} & ↑ \\
			S & Slovene\il{Slovene} & \cellcolor[RGB]{242,242,242} & \cellcolor[RGB]{229,229,229} & Serbian\il{Serbian} &  S \\
			& W & → &  → & E & \\
		\lspbottomrule
	\end{tabularx} \\

	\begin{tabularx}{0.75\textwidth}{lXXXXX}
		\lsptoprule
		\arrayrulecolor{customgray}
		& \textbf{Numeral} & \multicolumn{3}{l}{(only cardinal numeral)} & \\
		\cline{3-4}
		N & Polish\il{Polish} & \multicolumn{1}{|l|}{} & \multicolumn{1}{l|}{} & Russian\il{Russian} & N \\
		\cline{3-4}
		↓ & Czech\il{Czech} & \multicolumn{1}{|l|}{} & \cellcolor[RGB]{153,153,153} & Slovak\il{Slovak} & ↓ \\
		\cline{3-3}
		S & Slovene\il{Slovene} & \cellcolor[RGB]{153,153,153} & \cellcolor[RGB]{51,51,51} & Serbian\il{Serbian} &  S \\
		& W & → &  → & E & \\
		\arrayrulecolor{black}
		\lspbottomrule
	\end{tabularx}

	\begin{tabularx}{0.75\textwidth}{lXXXXX}
		\lsptoprule
		\arrayrulecolor{customgray}
		& \textbf{Adjective} & & & & \\
		\cline{4-4}
		N & Polish\il{Polish} & \cellcolor[RGB]{229,229,229} & \multicolumn{1}{l|}{} & Russian\il{Russian} & N \\
		↓ & Czech\il{Czech} & \cellcolor[RGB]{191,191,191} & \cellcolor[RGB]{191,191,191} & Slovak\il{Slovak} & ↓ \\
		S & Slovene\il{Slovene} & \cellcolor[RGB]{153,153,153} & \cellcolor[RGB]{153,153,153} & Serbian\il{Serbian} &  S \\
		& W & → &  → & E & \\
		\arrayrulecolor{black}
		\lspbottomrule
	\end{tabularx}
	\caption{Characterisation of \isi{uninflectability} in six Slavic languages}
	\label{03tab6}
\end{table}

In Table \ref{03tab6}, dark slots indicate a high degree of \isi{uninflectability}, while light slots mean a low degree. A white slot indicates total lack of \isi{uninflectability} for a certain part of speech, as is, e.g., the case for cardinal numerals in \ili{Russian}, \ili{Polish} and \ili{Czech} where all members of this class display overt inflectional forms. In \ili{Serbian}, on the other hand, numerals from `five' to `999' are all \isi{uninflectable}. A black slot thus means that a considerable number of lexical items is affected by \isi{uninflectability}.

The shading does not correspond to absolute numbers, but rather to a qualitative criterion: how many lexical classes are affected? How many types of words are affected by non-assignment and how regular is this? In \ili{Russian}, all classes of names, foreign words and abbreviations are affected by \isi{uninflectability}. In addition, nouns with any phonological ending and any \isi{gender} are affected, but systematically those ending in /e/, /i/ and /u/. Most varieties of abbreviations are also \isi{uninflectable}. Thus, \isi{uninflectability} of nouns also shows a high degree of regularity or systematicity. This is why the cell for the noun in Table \ref{03tab6} is shaded black for \ili{Russian}. In \ili{Slovene}, on the other hand, names or abbreviations are not affected by \isi{uninflectability} at all. Moreover, there is no phonological ending that would preclude nouns beforehand from being inflected. There are some \isi{uninflectable} nouns in \ili{Slovene}, but they cannot be described by rules. Therefore, the corresponding cell is shaded in light grey.

The conditions for this \isi{uninflectability} are different for the different parts of speech: In numerals we find a diachronic reduction of inflection over the past 100-150 years (\citealt{suprun1969}; \citealt[297--313]{doleschal2000}). In this case we are dealing with contextual uninflectedness\is{uninflectedness!contextual} as a step towards \isi{uninflectability}. Uninflectable\is{uninflectable} adjectives, on the other hand, are almost always borrowings (cf. \citealt[277--297]{doleschal2000}). Therefore, the reason for \isi{uninflectability} in adjectives can be found in their lack of integration and their foreignness. In the case of nouns, the situation is more complicated and will therefore be analysed in detail in the following section.

\is{uninflectability|(}
\section{Uninflectability in nouns} \label{03sec4}

\subsection{Introductory considerations}

Uninflectable\is{uninflectable} nouns can be divided into three lexically and morphologically defined classes, which overlap logically to some extent: \emph{proper nouns}, \emph{abbreviations} and \emph{borrowings}. These groups may intersect in the following way: many of the \isi{uninflectable} proper nouns are at the same time borrowings or foreign names, e.g., the (\ili{Italian}) male name \emph{Mario}, or the (\ili{German}) female name \emph{Sigrid} (in \ili{Russian}, \ili{Polish}, \ili{Czech}, \ili{Slovak}), the toponyms \emph{Oslo}, \emph{Cetinje} (in \ili{Russian}), and chrematonyms\footnote{Also called pragmatonyms, see \citet{zgusta1996}.} as \emph{Subaru} (in \ili{Russian}, \ili{Polish}, \ili{Czech}, \ili{Slovak}). It is very seldom the case that genuinely native names are \isi{uninflectable}, as, e.g., the \ili{Russian} male first name \emph{Sadko}\footnote{A mythical hero of the bylinas, cf. the opera by Rimsky-Korsakov.} or \ili{Russian} surnames as \emph{Dolgix}, \emph{Živago} in Russian and the \ili{Czech} type of surnames ending in /u:/ (graphic \textless ů\textgreater{}, such as \emph{Janů}, \emph{Petrů}) in \ili{Slovak}.\footnote{Both the \ili{Russian} and the \ili{Czech} surnames of these types are frozen Genitive forms. In \ili{Czech} such surnames may be inflected.} Abbreviations may of course also be borrowings and/or names (e.g., of organisations, such as \emph{IBM}), but they may also be native, such as \emph{SNG} `CIS' (\ili{Russian}), \emph{ORMO} (Ochotnicza Rezerwa Milicji Obywatelskiej `Volunteer Reserve of the Citizens' Militia') (\ili{Polish}) or compounds as \emph{kompolka} `commander of a regiment' (\ili{Russian}). The class of borrowings includes (besides names) all borrowed common nouns that are not morphologically integrated for various reasons, e.g. \emph{alibi}, \emph{tabu} (\ili{Russian}, \ili{Polish}, \ili{Czech}, \ili{Slovak}), \emph{miss} `miss', \emph{ledi} `lady' (all investigated languages).

Apart from these, there are small subgroups of native \isi{uninflectable} common nouns in some languages, e.g., female designations of persons (\emph{mami} `mummy' in \ili{Slovak} and \ili{Slovene}, \emph{professor} `woman-professor' etc. in \ili{Polish}),\footnote{The \ili{Polish} cases are, however, defective, in that they cannot be used in the plural *\emph{panie professor} `Ms:\textsc{n.pl} professor'.} extragrammatical\footnote{See \citet{dressler2000} for a discussion of this term.} formations, such as \emph{cucu} (\ili{Slovak}) or \emph{ciuciu} (\ili{Polish}) `candy', \emph{psipsi} (\ili{Polish}) `wee-wee' or the designations of letters and musical notes (\emph{be}, \emph{ce}, \emph{de} `b, c, d'; \emph{do} ``C'', \emph{re} ``D'', \emph{mi} ``E'' in all investigated languages besides \ili{Slovene}). Other native common nouns are very rarely \isi{uninflectable} (cf. e.g., \citealt[18]{krzyzanowski2013}).

In the literature on uninflectability various explanations have been brought forward for motivating the uninflectability of nouns. The main reason is that they do not fit the requirements of inflection, for phonological as well as morphological reasons (cf., e.g., \citealt[37]{krzyzanowski2013}). Thus, \citet[225]{isacenko1954} justifies the uninflectability of borrowings on /o/ in \ili{Russian}, which in the majority of cases bear final stress, with the lack of a morphonological model among \ili{Russian} nouns (the accent type \emph{okno} /oˈkno/ `window' comprises only a few \ili{Russian} native words and is not productive), and \citet{worth1966} attributes uninflectability in \ili{Russian} to a vocalic \isi{stem} ending, as would be the case in \emph{bo-a} `boa'. \citet[30]{krzyzanowski2013} mentions a word-final stressed vowel as one cause of uninflectability of abbreviations in \ili{Polish}, such as \emph{EKG} [ɛkaˈgɛ] `electrocardiogram'.

A notorious reason is also the mismatch between \isi{gender} and phonological shape of a noun, as with \emph{miss} `miss', \emph{madam} `madam' in \ili{Russian} \citep[148]{timberlake2004}. Another reason given is the recognisability of names, e.g., surnames that are homonymous with common nouns, such as \ili{Polish} \emph{Lato}, \emph{Musik} \citep{krzyzanowski2013} which is especially important in official style. This reason is also mentioned as a motivation for not inflecting \ili{Russian} place names as (neuter) \emph{Puškino}, which could be confused with the paronymous (masculine) place name \emph{Puškin} due to identical endings in the oblique cases (cf. \citealt{tejkin2021} for a discussion of this argument). Finally, foreignness is presented as a key factor for uninflectability, cf. the discussion by \citet[35--42]{krzyzanowski2013}, also \citet[1075--1076]{tejkin2021}.

It is also put forward that some nouns are assigned more easily or regularly to inflection than others. The latter either remain \isi{uninflectable} for a span of time or vary between inflection and \isi{uninflectedness}. This is, e.g., the case with borrowed common nouns and names ending in /o/ as well as with male names ending in /e/ or /i/ in \ili{Polish} (cf. \citealt[20--21]{krzyzanowski2013}).

But what does the fact that certain nouns do not (easily) fit into inflection imply for linguistic explanation? In this regard, two points can be made which will be explored in the next section:

\begin{enumerate}
	\item Systematic uninflectability demonstrates the limits and thereby the inner structure of the \isi{inflectional system} of a language. This is to say that in the case of systematic uninflectability there is a mismatch between some property of the noun in question and the rules conditioning inflection in a specific language.
	
	\item Unsystematic uninflectability is most often a phenomenon of prescriptive norms, but it may also reveal tendencies within inflection.
\end{enumerate}

\subsection{Uninflectability and inflectional classes\is{inflectional class|(}}

We shall commence with an examination of the phenomenon of unsystematic uninflectability. In \ili{Russian}, for instance, \ili{French} names ending in stressed /a/ are \isi{uninflectable} (as \emph{Nikol'a} [nʲɩkʌˈla] ``Nicolas'', \emph{Dega} [Dɛˈga] ``Degas''. This is not true for names in stressed /a/ from other languages (e.g., \emph{Abdulla} [ʌbdulˈla] ``Abdullah''). It is typical for unsystematic uninflectability that such a norm may change, e.g., \emph{Marseille} (ending in the soft consonant\is{consonant!soft} /j/) was \isi{uninflectable} in \ili{Slovak} according to the norm in 1998 and is now codified as both inflectable and \isi{uninflectable} (cf. \citealt{povazaj1998, povazaj2013}). Unsystematic uninflectability may therefore be indicative of a transitory state of a (set of) noun(s) towards inflection.

In other cases, unsystematic uninflectability may be an indicator of a tendency, as with nouns ending in /o/ in \ili{Polish}, where uninflectability is gaining ground both following and counter to prescriptive norms (cf. \citealt[56, 68--70, 98--100, 113, 126, 138]{doleschal2000} for discussion and references). E.g., male surnames in /o/ should be generally inflectable, but in dictionaries many foreign names are marked as \isi{uninflectable}, such as \emph{Eco}, \emph{Franco}. Other names in /o/, too, often remain \isi{uninflected} in usage (cf. \citealt[68--70]{doleschal2000}; \citealt[21, 41]{krzyzanowski2013}). Borrowed common nouns in /o/ with neuter \isi{gender}, on the other hand, are \isi{uninflectable} at first, but usually come to be inflected over time, e.g., \emph{agio} `agio' already in the 19\textsuperscript{th} century, \emph{auto} `car', \emph{radio} `radio', \emph{kino} `cinema' in the period before the Second World War, and more recently \emph{molo} `pier', \emph{kakao} `coacoa', which are classified as variably inflected and \isi{uninflected} (cf. \citealt[126]{doleschal2000}; \citealt{zmigrodzki2004-2021}; \citealt{wolinski2020}). The integration of neuters in /o/ into a corresponding inflectional class is therefore not automatic, but slower and not always complete (see \citealt[102]{dresslerdziubalskakolaczykfabiszak1997}; \citealt[112]{orzechowska1974}). This points to a development towards systematic uninflectability for nouns ending in /o/ in \ili{Polish}.

Let us now discuss systematic uninflectability as an indicator of the limits of the \isi{inflectional system}. Systematic uninflectability means that a certain type of noun cannot be assigned to any inflectional class of a given language and thus shows what the requirements of inflectional classes are.

The fundamental question is: what characteristics must a word possess, in order to be integrated into the \isi{inflectional system} of a language, i.e. to be assigned to an inflectional class? One important factor in the Slavic languages is the phonological ending. As previously stated, nouns ending in the vowel /u/ cannot be assigned to any inflectional class in \ili{Russian} or \ili{Polish}. This implies that the \ili{Russian} and \ili{Polish} inflectional classes do not permit a \isi{base form} in /u/. Consequently, inflectional classes can presuppose a certain \isi{base form} type, which is a conditio sine qua non for the assignment of a word to them. This type of property of an inflectional class will be referred to as a \textit{\isi{constitutive property}} (cf. \citealt{doleschal1997, doleschal2006}).

Usually native words correspond to the requirements of the \isi{inflectional system}, because lexicon and regular word formation of a given language provide words with a permissible \isi{base form}. Therefore, it is often foreign words that are affected by uninflectability, although this is not the only case: relict forms created with synchronously unproductive means of word formation can also be affected, e.g., the (archaic) \ili{Russian} first names \emph{Mixajlo}, \emph{Danilo} (cf. \citealt[739]{zaliznjak2008}). Therefore, the \isi{inflectional system} may also change in a way that leaves native words outside. Therefore, uninflectability enables us to identify, firstly, the constitutive properties\is{constitutive property} associated with inflectional classes and, secondly, the degree of productivity of the latter (cf. \citealt{wurzel1984, dressler1997, dressler2003}).

In the literature on Slavic languages, the integration of a noun into an inflectional class is usually related to \isi{gender} and to the phonological ending of the \isi{base form} as well as to the ending of the \isi{stem} \citep{hentschelmenzel2009}. In the case of the present Slavic languages, four properties have turned out to characterise inflectional classes of nouns: \emph{base form type}\is{base form}, \emph{stem type}\is{stem}, \emph{\isi{gender}} and \emph{semantic \isi{animacy}}.

As an example, let us consider first names. First names fall into two groups: male and female, which behave differently as to uninflectability. In the Slavic languages first names are characterised
by the following semantic and morpho(no)logical properties: they are inherently animate\is{animacy} and belong to either the masculine or the feminine \isi{gender}, where a canonical male first name belongs to the masculine \isi{gender} and a canonical female name to the feminine \isi{gender}. Depending on the respective language, their base forms\is{base form} may either end in a vowel, which in \ili{Czech}, \ili{Serbian}, \ili{Slovak} and \ili{Slovene} may be phonologically short or long (while in \ili{Russian} and Polish vowel length is not distinctive), or they may end in a consonant. Their stems\is{stem} may end in a vowel, a hard consonant\is{consonant!hard}, a soft consonant\is{consonant!soft} or a neutral\is{consonant!neutral} (\ili{Czech}, \ili{Slovak}) sc. unpaired\is{consonant!unpaired} (\ili{Russian}) consonant (cf. footnote \ref{03ftn2}).

It has now to be made clear, how a first name is assigned to an inflectional class, e.g., a borrowed name, such as the male name \emph{Pedro}. We can think of this assignment as a process of ``feature checking'': First \isi{animacy} and \isi{gender} are established -- in the case of \emph{Pedro}, masculine animate\is{animacy} \isi{gender}. Then the phonological ending of the name is matched with base forms\is{base form} of existing inflectional classes that contain masculine nouns. If one or several such classes exist, it has to be checked if the final vowel /o/ should be treated as part of the \isi{stem} or as the Nominative ending for this particular class. If the latter is the case, the word form is analysed as \emph{Pedr-o} `\textsc{stem}-ending'\is{stem} and it has to be checked in a final step if the resulting \isi{stem} conforms to the \isi{stem}-structure of the chosen inflectional class. If it does, the word becomes a member of this class. If it does not, it remains \isi{uninflectable}, unless it is adapted phonologically or morphologically, e.g., by changing the word ending and/or the \isi{stem} ending.

E.g., if the name \emph{Pedro} enters the \ili{Russian} language, it is assigned to masculine-animate\is{animacy} \isi{gender} (which means that it selects a specific class of agreement targets). Afterwards it is checked, if there is an inflectional class with a \isi{base form} ending in /o/ that allows for nouns with masculine-animate\is{animacy} \isi{gender}. It turns out that the only class with nouns ending in /o/ is reserved for neuter \isi{gender} (\emph{okno} \textsc{n.} `window'), whereas the classes allowing for masculine-animate\is{animacy} nouns
require a \isi{base form} ending in a consonant (\emph{Vladimir}) or in /a/ (\emph{Nikita}). Therefore, \emph{Pedro} cannot be assigned to any inflectional class and remains \isi{uninflectable} in \ili{Russian}.

If we take \ili{Slovene} on the opposite end of the noun continuum in Table \ref{03tab6}, the same process yields a different result: In \ili{Slovene}, there is one class specialized for masculine-animate\is{animacy} nouns which at the same time allows for a \isi{base form} in any vowel or consonant. In order to assign \emph{Pedro} to this class, the status of its \isi{stem} has to be checked. The declension in question allows both for stems\is{stem} ending in a consonant and in a vowel. Since the final /o/ in \emph{Pedro} is unstressed, it is analysed as a \textsc{nom.sg} ending and hence the noun is inflected, like the native name \emph{Marko}, as \emph{Pedr-o} (\textsc{nom}), \emph{Pedr-a} (\textsc{gen}), \emph{Pedr-u} (\textsc{dat}), \emph{Pedr-a} (\textsc{acc}), \emph{(o) Pedr-u} (\textsc{loc}), \emph{(s) Pedr-om} (\textsc{instr}).

Tables \ref{03tab7} and \ref{03tab8} show the degree of uninflectability in first names in the six languages under investigation. This degree is again marked by different colourings of the cells -- a darker colouring indicates a higher degree of uninflectability. A plus sign means that there are cases of systematic uninflectability for the \isi{base form} ending in question, while a plus in brackets indicates cases of unsystematic uninflectability for this cell. White colouring combined with a minus sign means total lack of uninflectability. Shaded cells with question marks indicate that such cases could not be verified and probably do not exist.

As can be deduced from Tables \ref{03tab7} and \ref{03tab8}, uninflectability is much more widespread with female first names in all Slavic languages under investigation.\footnote{In \ili{Polish} and \ili{Russian} vowel length is not distinctive, therefore the cells for the respective vowels are coalesced.} As to male first names, only \ili{Russian} and \ili{Polish} display cases of uninflectability, while in \ili{Slovak}, \ili{Czech}, \ili{Slovene} and \ili{Serbian} all male first names can be assigned to an inflectional class, regardless of their phonological make-up.

\begin{table}[ht]
	\centering
	\begin{tabularx}{0.83\textwidth}{cL{2cm}ccccccccccc}
		\lsptoprule
			\multirow{7}{*}{\rotatebox{90}{\hspace{-0.5cm}male first names}} & \textbf{base form ending in}\is{base form} & C\# & a & a: & o & o: & e & e: & i & i: & u & u: \\
			\cline{2-13}
			\addlinespace
			& Russian\il{Russian} & - & \multicolumn{2}{c}{\cellcolor[RGB]{191,191,191}+} & \multicolumn{2}{c}{\cellcolor[RGB]{64,64,64}\textcolor{white}{+}} & \multicolumn{2}{c}{\cellcolor[RGB]{0,0,0}\textcolor{white}{+}} & \multicolumn{2}{c}{\cellcolor[RGB]{0,0,0}\textcolor{white}{+}} & \multicolumn{2}{c}{\cellcolor[RGB]{0,0,0}\textcolor{white}{+}} \\
			& Polish\il{Polish} & - & \multicolumn{2}{c}{\cellcolor[RGB]{191,191,191}+} & \multicolumn{2}{c}{\cellcolor[RGB]{191,191,191}(+)} & \multicolumn{2}{c}{\cellcolor[RGB]{96,96,96}\textcolor{white}{(+)}} & \multicolumn{2}{c}{\cellcolor[RGB]{191,191,191}(+)} & \multicolumn{2}{c}{\cellcolor[RGB]{0,0,0}\textcolor{white}{+}} \\
			& Slovak\il{Slovak} & - & - & - & - & - & - & - & - & - & - & - \\
			& Czech\il{Czech} & - & - & - & - & - & - & - & - & - & - & - \\
			& Slovene\il{Slovene} & - & - & - & - & - & - & - & - & - & - & - \\
			& Serbian\il{Serbian} & - & - & - & - & - & - & - & - & - & - & - \\
		\lspbottomrule
	\end{tabularx}
	\caption{Uninflectability of male first names}
	\label{03tab7}
\end{table}


\begin{table}[ht]
	\centering
	\begin{tabularx}{0.92\textwidth}{cL{2cm}ccccccccccc}
		\lsptoprule
		\multirow{7}{*}{\rotatebox{90}{\hspace{-0.5cm}female first names}} & \textbf{base form ending in}\is{base form} & C\# & a & a: & o & o: & e & e: & i & i: & u & u: \\
		\cline{2-13}
		\addlinespace
		& Russian\il{Russian} & \cellcolor[RGB]{64,64,64}\textcolor{white}{+} & \multicolumn{2}{c}{\cellcolor[RGB]{217,217,217}+} & \multicolumn{2}{c}{\cellcolor[RGB]{0,0,0}\textcolor{white}{+}} & \multicolumn{2}{c}{\cellcolor[RGB]{0,0,0}\textcolor{white}{+}} & \multicolumn{2}{c}{\cellcolor[RGB]{0,0,0}\textcolor{white}{+}} & \multicolumn{2}{c}{\cellcolor[RGB]{0,0,0}\textcolor{white}{+}} \\
		& Polish\il{Polish} & \cellcolor[RGB]{0,0,0}\textcolor{white}{+} & \multicolumn{2}{c}{-} & \multicolumn{2}{c}{\cellcolor[RGB]{0,0,0}\textcolor{white}{+}} & \multicolumn{2}{c}{\cellcolor[RGB]{0,0,0}\textcolor{white}{+}} & \multicolumn{2}{c}{\cellcolor[RGB]{0,0,0}\textcolor{white}{+}} & \multicolumn{2}{c}{\cellcolor[RGB]{0,0,0}\textcolor{white}{+}} \\
		& Slovak\il{Slovak} & \cellcolor[RGB]{64,64,64}\textcolor{white}{+} & - & \cellcolor[RGB]{170,170,170}? & \cellcolor[RGB]{64,64,64}\textcolor{white}{+} & \cellcolor[RGB]{0,0,0}\textcolor{white}{+} & \cellcolor[RGB]{166,166,166}+ & \cellcolor[RGB]{0,0,0}\textcolor{white}{+} & \cellcolor[RGB]{0,0,0}\textcolor{white}{+} & \cellcolor[RGB]{0,0,0}\textcolor{white}{+} & \cellcolor[RGB]{0,0,0}\textcolor{white}{+} & \cellcolor[RGB]{0,0,0}\textcolor{white}{+} \\
		& Czech\il{Czech} & \cellcolor[RGB]{166,166,166}+ & - & \cellcolor[RGB]{170,170,170}? & \cellcolor[RGB]{166,166,166}(+) & \cellcolor[RGB]{166,166,166}(+) & \cellcolor[RGB]{217,217,217}(+) & \cellcolor[RGB]{166,166,166}(+) & \cellcolor[RGB]{0,0,0}\textcolor{white}{+} & - & \cellcolor[RGB]{0,0,0}\textcolor{white}{+} & \cellcolor[RGB]{0,0,0}\textcolor{white}{+} \\
		& Slovene\il{Slovene} & \cellcolor[RGB]{64,64,64}\textcolor{white}{(+)} & - & \cellcolor[RGB]{170,170,170}? & \cellcolor[RGB]{64,64,64}\textcolor{white}{(+)} & \cellcolor[RGB]{64,64,64}\textcolor{white}{(+)} & \cellcolor[RGB]{64,64,64}\textcolor{white}{(+)} & \cellcolor[RGB]{64,64,64}\textcolor{white}{(+)} & \cellcolor[RGB]{0,0,0}\textcolor{white}{+} & \cellcolor[RGB]{0,0,0}\textcolor{white}{+} & \cellcolor[RGB]{0,0,0}\textcolor{white}{+} & \cellcolor[RGB]{0,0,0}\textcolor{white}{+} \\
		& Serbian\il{Serbian} & \cellcolor[RGB]{127,127,127}\textcolor{white}{(+)} & - & \cellcolor[RGB]{170,170,170}? & \cellcolor[RGB]{96,96,96}\textcolor{white}{(+)} & \cellcolor[RGB]{96,96,96}\textcolor{white}{(+)} & \cellcolor[RGB]{127,127,127}\textcolor{white}{(+)} & \cellcolor[RGB]{96,96,96}\textcolor{white}{(+)} & \cellcolor[RGB]{0,0,0}\textcolor{white}{+} & \cellcolor[RGB]{0,0,0}\textcolor{white}{+} & \cellcolor[RGB]{0,0,0}\textcolor{white}{+} & \cellcolor[RGB]{0,0,0}\textcolor{white}{+} \\
		\lspbottomrule
	\end{tabularx}
	\caption{Uninflectability of female first names}
	\label{03tab8}
\end{table}

We will now concentrate on \ili{Russian} first names and compare them to \ili{Slovene} first names.

\ili{Russian} inflectional classes are traditionally divided into three macroclasses\footnote{According to \citet{dressler1997, dressler2003} a class is a set of similar paradigms\is{paradigm} that can be subdivided into subclasses and	ultimately into ``microclasses which share exactly the same morphological generalizations. (\ldots) An inflectional macroclass is the highest, most general type of class, which comprises several	(sub)classes (\ldots) Prototypically, its nucleus is a productive microclass and it has at least two microclasses. The interior coherence of a macroclass, in terms of shared properties, must be higher than affinities between microclasses of different macroclasses'' \citep[35--36]{dressler2003}.} according to common sets of endings (cf. for \ili{Russian}, e.g., \citealt[130]{timberlake2004}). The 1\textsuperscript{st} macroclass displays (among others) a class containing nouns with masculine animate\is{animacy} \isi{gender} and a \isi{base form} ending in a consonant, which is divided further into subclasses according to the \isi{stem} ending (identical to the \isi{base form} ending) and \isi{animacy}. Typical \ili{Russian} male names belonging to this class are, e.g., \emph{Boris} (ending in a hard consonant\is{consonant!hard}), \emph{Igor'} (ending in a soft consonant\is{consonant!soft}), \emph{Jurij} (ending in an unpaired consonant\is{consonant!unpaired}). The 2\textsuperscript{nd} macroclass, on the other hand, requires a \isi{base form} ending in /a/ and a \isi{stem} ending in a consonant and can again be divided into subclasses depending on \isi{animacy}. Animate\is{animacy} nouns in this class may have feminine or masculine \isi{gender}. Typical \ili{Russian} female names here are, e.g., \emph{Anna} (\isi{stem} ending in a hard consonant\is{consonant!hard}), \emph{Marija} (stem ending in an unpaired consonant\is{consonant!unpaired}), (hypocoristic) \emph{Anja} (\isi{stem} ending in a soft consonant\is{consonant!soft}), typical male names are \emph{Nikita} (\isi{stem} ending in a hard consonant\is{consonant!hard}) and (hypocoristic) \emph{Kolja} (\isi{stem} ending in a soft consonant\is{consonant!soft}). The 3\textsuperscript{rd} macroclass has a subclass with feminine animate\is{animacy} \isi{gender} and a \isi{base form} (as well as \isi{stem}) ending in a soft consonant\is{consonant!soft} which contains also traditional \ili{Russian} female names, e.g., \emph{Ljubov'}, \emph{Raxil'} ``Rachel''.

Let us first look at male names. As can be seen from Table \ref{03tab7}, in \ili{Russian} male names ending in a vowel are affected by systematic uninflectability. But only in case of the vowels /e/, /i/ and /u/ is this correlation 100\%, i.e., it is the \isi{base form} ending alone that conditions uninflectability. In other words, masculine-animate\is{animacy} proper nouns with a \isi{base form} ending in /e/, /i/ or /u/, do not conform to the conditions of any inflectional class in \ili{Russian}. Therefore, foreign names as \emph{Arne, Teddi, Lu} are \isi{uninflectable}:

\begin{exe}
	\ex \gll Ja vstretila Arne/Teddi/Lu. \\
	I met:\textsc{f} Arne.\textsc{acc}/Teddy.\textsc{acc}/Lou.\textsc{acc} \\
	\trans `I met Arne/Teddy/Lou.'
\end{exe}

With names ending in /o/ and /a/, other factors come into play, as can be seen in Table \ref{03tab9}. Names ending in /o/ are \isi{uninflectable}, if the /o/ is stressed (\emph{Sadko}):

\begin{exe}
	\ex \gll Ja vstretila Sadko. \\
	I met:\textsc{f} Sadko.\textsc{acc} \\
	\trans `I met Sadko.'
\end{exe}

Nevertheless, some traditional Slavic names ending in unstressed /o/ (\emph{Miхajlo})\footnote{This form of the name ``Michael'' (nowadays \emph{Mixail}) was the original given name of the great \ili{Russian} scholar Mixail Lomonosov. Such forms (as also \emph{Danilo}, \emph{Ivanko}) are obsolete as given names today. But they are still presented in dictionaries \citep{petrovskij-nodate} and characterised as inflectable according to the 2\textsuperscript{nd} declension which normally requires a \isi{base form} ending in /a/, see \citeauthor{zaliznjak2008} (\citeyear[739]{zaliznjak2008}, also on \emph{Sadko}).} may be inflected according to the 2\textsuperscript{nd} inflectional macroclass:

\begin{exe}
	\ex \gll Ja vstretila Mixajl-u. \\
	I met:\textsc{f} Mixajlo-\textsc{acc} \\
	\trans `I met Mixajlo.'
\end{exe}

This fact is reflected in Table \ref{03tab9} by putting this \isi{base form} in brackets. (Unstressed /o/ is pronounced like unstressed /a/ as [ə], which may partly motivate this class assignment).

In male names ending in /a/, uninflectability is conditioned by the \isi{stem} ending. The corresponding nouns are assigned to the 2\textsuperscript{nd} declension, if /a/ is analysed as the Nominative suffix -a, as in \emph{Nikit-a}):

\begin{exe}
	\ex \gll Ja vstretila Nikit-u. \\
	I met:\textsc{f} Nikita-\textsc{acc} \\
	\trans `I met Nikita.'
\end{exe}

But membership in this class is allowed only if the remaining \isi{stem} ends in a consonant. Therefore, a name such as (\ili{French}) \emph{Fransua} ``François'' whose \isi{stem} would end in /u/ (*Fransu-a) cannot be inflected:

\begin{exe}
	\ex \gll Ja vstretila Fransua. \\
	I met:\textsc{f} François.\textsc{acc} \\
	\trans `I met François.'
\end{exe}

But (Japanese) \emph{Akira} is inflectable, since its \isi{stem} ends in the consonant /r/ (Akir-a):

\begin{exe}
	\ex \gll Ja vstretila Akir-u. \\
	I met:\textsc{f} Akira-\textsc{acc} \\
	\trans `I met Akira.'
\end{exe}

Hence in \ili{Russian} only male names ending in a consonant can be integrated into inflectional classes unconditionally, because in a concrete noun, all four constitutive properties\is{constitutive property} -- \isi{base form} type, \isi{stem} type, \isi{gender}, subgender -- have to have the fitting value. This implies that the constitutive properties\is{constitutive property} of the \ili{Russian} nominal inflectional classes are applied in a restrictive way and thus motivates the high degree of systematic uninflectability in male names.

\begin{table}[ht]
	\centering
	\begin{tabularx}{\textwidth}{>{\raggedright}p{1.3cm}p{1cm}p{1.4cm}p{1.8cm}p{1.6cm}Q}%L{1.3cm}p{1cm}p{1.4cm}p{1.8cm}p{1.6cm}L{3cm}
		\lsptoprule
			Class 					   		   & Base form\is{base form} & Gender\is{gender} & Stem type\is{stem} & Inflectable example & Uninflectable\is{uninflectable} example \\
			\midrule
			1. hard\is{consonant!hard} 		   & C-Ø 					 & \textsc{m-anim.}  & hard C 			  & Boris-Ø, Bob-Ø & Pedro, Arne, Lu “Lou” \\
			1. soft\is{consonant!soft} 		   & C-Ø 					 & \textsc{m-anim.}  & soft C 			  & Igor'-Ø & Mikele ``Michele'', Teddi ``Teddy'' \\
			1. unpaired\is{consonant!unpaired} & C-Ø 		     		 & \textsc{m-anim.}  & unpaired C 		  & Jurij-Ø & Mat'je ``Mathieu'', Luidži ``Luigi'' \\
			2. hard 				   		   & -a   				     & \textsc{m-anim.}  & hard C 			  & Nikit-a & Fransua ``François'' \\
									   		   & (-o) 					 & 					 & 					  & (Mixajl-o) & Sadko (stressed /o/) \\
			2. soft 				   		   & -a 				     & \textsc{m-anim.}  & soft C 			  & Kolj-a & Nikolja ``Nicolas'' (stressed /a/ and French\il{French} origin) \\
		\lspbottomrule
	\end{tabularx}
	\caption{Constitutive properties\is{constitutive property} of \ili{Russian} masculine inflectional classes\protect\footnotemark}
	\label{03tab9}
\end{table}

\footnotetext{In \ili{Russian} /j/ belongs to the “unpaired consonants\is{consonant!unpaired}”.}

The inflectional classes of \ili{Slovene} are essentially analogous to those of \ili{Russian}. Consequently, we will adhere to the numerical classification previously outlined.\footnote{Note that \ili{Slovene} grammaticography defines macroclasses on the basis of \isi{gender}, cf. \citet[40--61]{herrity2000}; \citet[28--31]{greenberg2008}.} The 1\textsuperscript{st} macroclass displays again a class for nouns with masculine-animate\is{animacy} \isi{gender}. The 2\textsuperscript{nd} macroclass is characterised by a \isi{base form} in /a/ and includes one class comprising feminine and masculine-animate\is{animacy} nouns. (The 3\textsuperscript{rd} macroclass is reserved for feminine \isi{gender} nouns, and is not productive in integrating first names).

With regard to male first names, these can belong to either the 1\textsuperscript{st} or the 2\textsuperscript{nd} macroclass (Table \ref{03tab10}). In the 1\textsuperscript{st} class, the \isi{base form} of native nouns typically ends in a consonant (e.g., \emph{Borut, Matjaž} ``Matthew''), in the 2\textsuperscript{nd}, it typically ends in /a/ (as in \emph{Miha}\footnote{Male names in unstressed /a/ can inflect either according to the 2\textsuperscript{nd} or the 1\textsuperscript{st} macroclass.} ``Michael''). Stems\is{stem} typically end in a consonant, as in \ili{Russian}. However, in contrast to \ili{Russian}, these inflectional classes permit a wide range of base forms\is{base form} and also allow for vocalic stems\is{stem}, as evidenced by the \ili{Latin} name \emph{Leo}, where the final, unstressed /o/ is analysed as an ending and the noun inflected accordingly as \emph{Le-o} (\textsc{nom}), \emph{Le-a} (\textsc{gen}), \emph{Le-u} (\textsc{dat}), \emph{Le-a} (\textsc{acc}), \emph{(o) Le-u} (\textsc{loc}), \emph{(s) Le-om} (\textsc{instr}). If the name ends in a stressed vowel or in /e/, /i/, /u/, the \isi{stem} is automatically augmented by /j/\footnote{In \ili{Slovene} /j/ belongs to the ``soft consonants\is{consonant!soft}''.} (or sometimes /t/) in the oblique cases, thereby creating a consonantal \isi{stem} ending. This can be observed in the examples of \emph{Teddy-Ø, Teddyj-a} (\textsc{gen}), \emph{Lou-Ø, Louj-a} (\textsc{gen}), \emph{Tone-Ø}, \emph{Tonet-a} (\textsc{gen}) ``Tony''. Consequently, all male names are automatically assigned to an inflectional class. Therefore there are no \isi{uninflectable} male names in \ili{Slovene}.

\begin{table}[ht]
	\centering
	\begin{tabularx}{\textwidth}{lp{1.6cm}lQQl}%lp{1.6cm}p{1.4cm}L{2cm}L{3cm}l
		\lsptoprule
		Class & Base form\is{base form} & Gender\is{gender} & Stem type\is{stem} & Inflectable & Uninfl. \\
		\midrule
		1. hard\is{consonant!hard} & C-Ø, \newline -o, -a, -e \newline (unstressed) & \textsc{m-anim.} & hard C or V & Borut-Ø, Mark-o, Mih-a \newline Bob-Ø, Pedr-o, Le-o, Akir-a & - \\
		\addlinespace
		1. soft\is{consonant!soft} & C-Ø \newline V-Ø & \textsc{m-anim.} & soft C (\isi{stem}-augmentation after vowel) & Matjaž-Ø, Franz-Ø \newline Nicolas-Ø (stressed /a/), François-Ø (stressed /a/) \newline Arne-Ø, Teddy-Ø, Lou-Ø & - \\
		\addlinespace
		2. & -a & \textsc{m-anim.} & C & Mih-a, Akir-a & - \\
		\lspbottomrule
	\end{tabularx}
	\caption{Constitutive properties\is{constitutive property} of \ili{Slovene} masculine inflectional classes}
	\label{03tab10}
\end{table}

We will now consider female first names. As previously noted, female names are significantly affected by uninflectability, as illustrated in Table \ref{03tab8}. In \ili{Russian}, all female names ending in a vowel other than /a/ are systematically \isi{uninflectable}, e.g., \emph{Renate, Mimi, Ildiko, Lulu}. Furthermore, names with a \isi{base form} in a consonant are also characterised by a high degree of uninflectability, e.g., \emph{Èdit} (\ili{German}) ``Edith'', \emph{Kler} ``Claire'', and even certain names with a \isi{base form} in /a/ may be \isi{uninflectable}, such as \emph{Klaudia} ``Claudia'', \emph{Rea} ``Rhea'' (with two coda vowels).

In \ili{Slovene}, uninflectability is also found in female first names, but systematically only with names ending in /i/ and /u/, such as \emph{Mimi}, \emph{Lulu}. Similarly, names with a \isi{base form} ending in the vowels /e/ and /o/ as well as in a consonant are also affected by uninflectability, although not in a systematic manner: a corpus search for \emph{Renate Götschl}\footnote{A well-known Austrian athlete (1990's-2000's).} in the \ili{Slovene} corpus \citetitle{gigafida2024} yields 675 tokens for the word-form \emph{Renate}, and one token for the inflected form \emph{Renat-i} (\textsc{dat}) (excluding items that may be related to a \isi{base form} \emph{Renata}). The names \emph{Isabel} or \emph{Madeleine} usually appear in an \isi{uninflected} form: A total of 3,123 tokens of the form \emph{Madeleine} are observed, along with one instance of the form \emph{Madelein-o} (\textsc{acc}) and two instances of the form \emph{Madelein-i} (\textsc{dat}). The same is true, mutatis mutandis, for \emph{Isabel}: there are 1,926 tokens of \emph{Isabel} and two of \emph{Isabel-i} (\textsc{dat}) (where the relation to this \isi{base form}, rather than an adapted \emph{Isabela} is undisputable). Names ending in consonant or /e/ are therefore classified as unsystematically \isi{uninflectable}. In contrast, names ending in /a/ are unconditionally inflectable and go into the 2\textsuperscript{nd} (macro)class.

\largerpage[2]
Let us again turn our attention to the inflectional classes of \ili{Russian} and \ili{Slovene}. In \ili{Russian}, the conditions for female first names are similar to those for male first names: in order to be assigned to an inflectional class, a female name with feminine-animate\is{animacy} \isi{gender} has to provide a \isi{stem} ending in a consonant and a \isi{base form} ending in either /a/ or a soft consonant\is{consonant!soft} (the 3\textsuperscript{rd} macroclass). In reality, however, this latter case is obsolete, since the third declension is no longer productive (in the sense of \citealt{dressler1997, dressler2003}). This means that it is not open for borrowings or extragrammatical formations, even if they fulfil the conditions of the constitutive properties\is{constitutive property}, e.g., the Chinese name \emph{Junin'} ``Yuning''. Nevertheless, older borrowings, such as biblical names (e.g., \emph{Raxil'} ``Rachel''), or the Soviet extragrammatical formation \emph{Ninel'} (an anagram of the name \emph{Lenin}) are inflectable (where \emph{Ninel'} often remains \isi{uninflected}, as a search of the \citetitle{savcuketal2024}\footnote{The research was conducted using the \citetitle{savcuketal2024} (\href{https://ruscorpora.ru/}{ruscorpora.ru}), \citet{savcuketal2024}.}
confirms). Therefore, this type is not classified as systematically \isi{uninflectable}. New feminine\footnote{This \isi{gender} feature has to be emphasized, because in \ili{Russian} there are hypocoristic female names with masculine \isi{gender}, which are derived by a native suffix and are inflected according to the 1\textsuperscript{st} macroclass, e.g.: \emph{Liza} \textsc{f.} \textgreater{} \emph{Liz-oček} \textsc{m.}} female names can only be integrated into inflection if they have a \isi{base form} that ends in /a/ and a \isi{stem} that ends in a consonant. Therefore, names as \emph{Rea} (*Re-a), \emph{Klaudia} (*Klaudi-a) remain \isi{uninflectable}.

Female names can of course be adapted by changing an unfitting \isi{base form} ending into /a/ (e.g., \emph{Renate} \textgreater{} \emph{Renata}) or by inserting the glide /j/ between two coda vowels (e.g., \emph{Klaudia} \textgreater{} \emph{Klaudija}).\footnote{\citet{ageenko2010} lists \emph{Klaudija Kardinale} (with glide-insertion) and \emph{Klaudia Schiffer} (without), \citet{klysinskij2010} recommends the form with /j/ for the transcription of \ili{Italian} names in \textless ia\textgreater{} but without /j/ for \ili{German} names with this ending.} (This occurs frequently in casual speech and in the dialects). However, this is not a common practice in Contemporary Standard \ili{Russian}, where usually unfitting names remain as they are and are therefore \isi{uninflectable}.

\begin{table}[ht]
	\centering
	\begin{tabularx}{\textwidth}{lQp{1.3cm}QlQ}%lp{1.3cm}p{1.3cm}L{2.2cm}lL{2cm}
		\lsptoprule
		Macroclass & Base form\is{base form} & Gender\is{gender} & Stem type\is{stem} & Inflectable & Uninfl. \\
		\midrule
			2. hard\is{consonant!hard} & -a & \textsc{f-anim.} & hard C & Anna & Rea ``Rhea'' \\
			2. soft\is{consonant!soft} & -a & \textsc{f-anim.} & soft C & Anja & \\
			3. & C-Ø & \textsc{f-anim.} & soft C, unpaired\is{consonant!unpaired} /š/, /ž/ & Raxil' & Solanž ``Solange'', Junin' ``Yuning'' \\
		\lspbottomrule
	\end{tabularx}
	\caption{Constitutive properties\is{constitutive property} of \ili{Russian} feminine inflectional classes}
	\label{03tab11}
\end{table}

In \ili{Slovene} the situation is more fluid, although the tendency is essentially the same: only female names ending in /a/ are fully inflectable and are assigned to one of the subclasses of the 2\textsuperscript{nd} macroclass. In contrast to \ili{Russian}, this class displays no restrictions as to the properties of \isi{animacy} and \isi{stem} type. Therefore, both animate\is{animacy} and inanimate\is{animacy} nouns with a consonantal or vocalic \isi{stem} can be assigned to it. Furthermore, names with a \isi{base form} ending in /o/, /e/ or a consonant can be assigned to this class. The part before such a vowel is analysed as the \isi{stem} to which inflectional suffixes are attached. However, such cases are very rare and often historical, as with names from Greek mythology or from the Bible (\emph{Kli-o} (\textsc{nom}), \emph{Kli-e} (\textsc{gen})). The 3\textsuperscript{rd} macroclass of \ili{Slovene}, finally, does not contain or integrate female names at all.

\begin{table}[ht]
	\centering
	\begin{tabularx}{\textwidth}{lp{1.3cm}lp{0.8cm}QQ}%lL{1.8cm}p{1.2cm}p{1.3cm}L{2.8cm}L{2cm}
		\lsptoprule
		Class & Base form\is{base form} & Gender\is{gender} & Stem type\is{stem} & Inflectable & Uninflectable\is{uninflectable} \\
		\midrule
			2. & -a, (o, e, C-Ø) & f & any & Alina, Lea, (Klio, Melpomene, Ruth, Madeleine, Isabel) & Lulu, Mimi, Ildiko \\
			3. & C-Ø & f & C & no names & no names \\
		\lspbottomrule
	\end{tabularx}
	\caption{Constitutive properties\is{constitutive property} of \ili{Slovene} feminine inflectional classes}
	\label{03tab12}
\end{table}

This discussion has shown that the nominal inflectional classes of \ili{Russian} and \ili{Slovene} are characterised by the constitutive properties\is{constitutive property} \emph{\isi{gender}} (including the subgender of \emph{\isi{animacy}}), \emph{base form type}\is{base form} and \emph{\isi{stem} type}. In these cases, the semantic category of \emph{\isi{animacy}} is grammaticalised both as a subgender and a subclass of the corresponding inflectional classes.

However, \isi{animacy} may also be a semantic \isi{constitutive property}, as evidenced by \ili{Czech} and \ili{Slovak}. In \ili{Czech}, there is one inflectional class for feminine nouns ending in a soft consonant\is{consonant!soft}, analogous to the respective 3\textsuperscript{rd} class in \ili{Russian} and \ili{Slovene}. This class is productive, but it does not contain animate\is{animacy} nouns. Because of this semantic \isi{constitutive property}, it integrates common nouns such as \emph{kartridž-Ø} `cartridge', toponyms as \emph{Marseille} (ending in /j/) and \emph{Provence} (ending in /s/), but a female name such as \emph{Florence} (ending in /s/) cannot be assigned to this class and remains \isi{uninflectable}\is{uninflectability|)}\is{inflectional class|)}.

\section{Summary}

The objective of this article is to demonstrate the following points: Uninflectedness\is{uninflectedness} in inflecting languages is not a uniform phenomenon. In the first place, two fundamentally distinct manifestations can be identified: contextual\is{uninflectedness!contextual} and paradigmatic uninflectedness\is{uninflectedness!paradigmatic}. The former represents a temporary absence of inflection in an otherwise inflectable lexeme, whereas the latter is a permanent feature of either a single lexeme or a class of lexemes with identical properties. The first type has been termed \emph{contextual uninflectedness}\is{uninflectedness!contextual}, while the two varieties of the second type have been termed \emph{systematic} and \emph{unsystematic \isi{uninflectability}}. These manifestations can be attributed to different causes which have been spelled out for the noun.

Contextual uninflectedness\is{uninflectedness!contextual} is caused by contextual factors, such as syntactic context or register, but may also be sociolinguistically conditioned. Systematic \isi{uninflectability} of nouns is caused either by a mismatch between the constitutive properties\is{constitutive property} of inflectional classes\is{inflectional class} and the properties of the nouns in question, or by the fact that an otherwise matching class\is{inflectional class} is unproductive. Finally, unsystematic \isi{uninflectability} can on the one hand occur for reasons of style or register. On the other hand, it can be an indicator of ongoing changes, showing either that an \isi{inflectional class} is changing its constitutive properties\is{constitutive property} or that it is becoming unproductive\il{Slavic|)}.
 
 
%\section*{Abbreviations}
%\begin{tabularx}{.45\textwidth}{lQ}
%... & \\
%... & \\
%\end{tabularx}
%\begin{tabularx}{.45\textwidth}{lQ}
%... & \\
%... & \\
%\end{tabularx}

\section*{Acknowledgements}
I am very much obliged to my reviewers for their comments which have significantly contributed to improving this chapter. All flaws are of course mine.

%\section*{Contributions}
%John Doe contributed to conceptualization, methodology, and validation.
%Jane Doe contributed to writing of the original draft, review, and editing.


{\sloppy\printbibliography[heading=subbibliography,notkeyword=this]}
\end{document}
